\documentclass[10pt]{amsart}
\usepackage{amsmath,amssymb,amsthm}
% US letter, margins of about 1.8 cm set by hand; compact section headings.
\setlength{\textwidth}{\paperwidth}\addtolength{\textwidth}{-3.6cm}
\setlength{\oddsidemargin}{-0.74cm}\setlength{\evensidemargin}{-0.74cm}
\setlength{\headheight}{0pt}\setlength{\headsep}{0pt}
\setlength{\topmargin}{-0.84cm}
\setlength{\textheight}{\paperheight}\addtolength{\textheight}{-3.6cm}
\setlength{\footskip}{0.8cm}
\ifdefined\pdfpagewidth\pdfpagewidth=\paperwidth\pdfpageheight=\paperheight\fi
\pagestyle{plain}
\setlength{\parskip}{0pt plus 0.5pt}
\widowpenalty=150 \clubpenalty=150
\makeatletter\def\section{\@startsection{section}{1}{\z@}{.6\baselineskip}{.25\baselineskip}{\normalfont\bfseries}}\makeatother

\numberwithin{equation}{section}
\newtheoremstyle{compact}{4pt}{4pt}{\itshape}{}{\bfseries}{.}{.5em}{}
\theoremstyle{compact}
\newtheorem{theorem}{Theorem}[section]
\makeatletter % the proof environment of amsthm, with the spacing of the theorems
\renewenvironment{proof}[1][\proofname]{\par\pushQED{\qed}\normalfont
  \topsep4\p@\relax\trivlist\item[\hskip\labelsep\itshape#1\@addpunct{.}]%
  \ignorespaces}{\popQED\endtrivlist\@endpefalse}
\makeatother

\newcommand{\CP}{\mathbb{CP}}
\newcommand{\CPb}{\overline{\mathbb{CP}}{}^2}
\newcommand{\Z}{\mathbb Z}
\newcommand{\PD}{\operatorname{PD}}
\newcommand{\Sym}{\operatorname{Sym}}
\newcommand{\SO}{\mathrm{SO}}
\newcommand{\fs}{\mathfrak s}
\newcommand{\ft}{\mathfrak t}
\newcommand{\cZ}{\mathcal Z}
\newcommand{\cV}{\mathcal V}
\newcommand{\cM}{\mathcal M}
\newcommand{\step}[1]{\par\noindent\textit{#1}\enspace}
% References set as one paragraph, to keep the note within two pages.
\makeatletter
\renewenvironment{thebibliography}[1]{\par\footnotesize\noindent
  \textbf{References.}\ignorespaces
  \def\@lbibitem[##1]##2{\unskip\enspace[##1]~\if@filesw\immediate
    \write\@auxout{\string\bibcite{##2}{##1}}\fi\ignorespaces}%
  \def\bysame{---}}{\par}
\makeatother

\begin{document}

\section{Strategy}

Let $K\subset S^3$ be a nontrivial knot, $Y_r=S^3_r(K)$, and
$\phi\colon Y_{-2}\to Y_2$ an orientation-preserving diffeomorphism. By
\cite[Cor.~1.4]{DLME}, after Ni--Wu and Hanselman, this is the last open case
of the cosmetic surgery conjecture, and the Chern--Simons filtration, which
settles $\pm1/n$, is unfavourable for the analogous maps, such as $B$ below
\cite[\S1.2]{DLME}. We replace it by an obstruction from $\SO(3)$-monopoles. Thanks to $\phi$,
copies of $B$ compose indefinitely, keeping the pairings of
Step~2 nonzero, but each brings a sphere of square $-4$ whose $\mu$-class
costs every Seiberg--Witten stratum more energy than it supplies (Step~3),
until no stratum can carry the nonzero count that Step~1 requires.

\step{Step 1. A Witten-type relation, used in reverse.}
The $\SO(3)$-monopole cobordism of Pidstrigach and Tyurin, from which Feehan
and Leness derive Witten's conjecture \cite{W} for many four-manifolds
modulo a gluing theorem \cite{PT,FLM,FL6}, also gives a proved vanishing
theorem. Take a spin$^u$ structure $\ft$ (a spin$^c$ structure
twisted by a $\mathrm U(2)$-bundle $E$) with $w=c_1(E)$, $\Lambda=c_1(\ft)$,
energy $\kappa=-p_1(\ft)/4$, and positive index
$n_D=(\Lambda^2-\sigma)/4-\kappa$ of the Dirac operator coupled to an
instanton. Near a generic instanton $A$ the link in the quotient by the circle
acting on the spinor has fibre $\mathbb P(\ker D_A)\cong\CP^{n_D-1}$, on which the
class $\mu_c$ of \cite{FL2b} is twice the hyperplane class. So for $\deg
z=\dim M^w_\kappa$ the quotient cut down by $\mu_p(z)$ and $\mu_c^{n_D-1}$ is a
one-manifold with $2^{n_D-1}D^w_X(z)$ ends at instantons
\cite[Prop.~3.29]{FL2b}; its other ends are at Seiberg--Witten strata
$M_\fs\times\Sym^\ell(X)$, $c_1(\fs)=\Lambda+v$, of level
$\ell=\kappa+v^2/4\ge0$, the number of bubbles. If it reaches none,
Stokes' theorem gives $D^w_X(z)=0$ without gluing at reducibles
(Theorem~\ref{thm:closed}). In reverse: a nonzero count must reach a
stratum with $M_\fs\neq\emptyset$, even if $SW(\fs)=0$; we produce one where
none can be reached.

\step{Step 2. Nonvanishing.}
In Floer's instanton homology $I$ \cite{DLME}, let $B\colon I(Y_2)\to
I(Y_{-2})$ be the map of the negative-definite cobordism
$W'=(-X_2(K))\cup_{S^3}X_{-2}(K)$ over an interval of metrics (Step~4), and
$H\colon I(Y_{-2})\to I(Y_2)$ that of the cobordism $W_H$ through $Y_0$. Both
are defined over $\Z$; $H$ is an isomorphism, and modulo two $B$ is expected to invert
it up to a fixed automorphism. So $u=(\phi_*B)^3HB$ is invertible on
$(I(Y_2;\Z)/\text{torsion})\otimes\Z_{(2)}$, and $u^M\equiv1\pmod{2^N}$ for
some $M$. Caps of $\pm Y_2$ with $b^+>0$ and relative
invariants $\langle\Psi_+,\Psi_-\rangle=q\neq0$, obtained as in \cite{KM04}
from a symplectic manifold containing $Y_0$ (here $K$ must be nontrivial),
give
$\langle\Psi_+,u^M\Psi_-\rangle\equiv q\not\equiv0\pmod{2^N}$. Integrality
matters: Step~1 has a factor $2^{n_D-1}$, which vanishes modulo two.

\step{Step 3. An insertion costs more energy than it supplies.}
The cores of the $2$-handles of $W'$ form a sphere $S$, $[S]=F_l-F_r$,
with halves $F_l$, $F_r$ from the two traces, orthogonal of square $-2$; so
$S\cdot S=-4$, $\partial\nu S=L(4,1)$, and $B$ is cut down by $\mu(S)$. Regard
$\kappa$ as the energy available; a stratum spends $-v^2/4$ on its
reducible and $\ell$ on bubbles. As
$\dim M^w_\kappa=8\kappa-3(1+b^+)$, a degree-two insertion supplies $1/4$, or
$1/8$ with a parameter as in $B$. Represent $\mu(S)$ by the
jumping-line divisor $V_S$ (\S\ref{sec:closed}). A reducible with class
$v$ lies in $V_S$ iff $\langle v,S\rangle\neq0$, and then the halves,
on which $w$, hence $v$, is even, cost at least $1/2$ in $-v^2/4$; if
$\langle v,S\rangle=0$, a limit lies in $V_S$ only with a bubble on $S$. So a
stratum is reached only if $\ell\ge T(v)$, the number of spheres orthogonal
to $v$, and each sphere lowers $\ell-T(v)$ by at least
$3/8$.

\step{Step 4. Why a family.}
The closed manifold $X_M$ capping the composite of $u^M$ is a
connected sum along the middle $3$-spheres of the copies of $W'$, with
$b^+>0$ on both sides, so its Donaldson and Seiberg--Witten invariants vanish.
The pairing of Step~2 computes instead a secondary invariant $\Omega_M$, a
count of instantons over the cube $[-1,1]^n$ of metrics whose faces stretch the necks $S^3$
and $L(4,1)$, cut down by the $\mu(S_i)$ and the cap classes, and summed with
signs over the bundles $c_e=c_0+\sum_ie_i\PD(S_i)$, $e\in\{0,1\}^n$. So is $B$:
summed over $c$ and $c+\PD(S)$, the map of $W'$ cut down by $\mu(S)$ is
null-homotopic by stretching $S^3$, as $I(S^3)=0$, and $L(4,1)$, where the
ends at the reducible of energy $1/4$ on $\nu S$ cancel; $B$ is the
difference. The
parameter also lowers the energy an insertion supplies from $1/4$ to
$1/8$, and this is decisive: $\ell(K)+\ell(-K)-2n_D=4\kappa+(K^2+\sigma)/2$ does not involve $\Lambda$, and a sphere with its halves changes
it by $4s-1$, $s$ the energy supplied, for $K$ orthogonal to the halves (the
hardest to exclude): by $0$ in a closed manifold, by $-1/2$ in a family. So with
$\Lambda$ constrained as in Step~5, a closed manifold gains nothing from
spheres (Section~\ref{sec:closed}).

\step{Step 5. Why $b^+$.}
Step~1 needs $n_D\ge1$, else the link is empty; Feehan and Leness
take $\Lambda$ large \cite[Thm.~1]{FLM}. Here the relation is applied to the signed sum
over the $c_e$, whose instanton ends carry $2^{n_D-1}$, so $n_D$ must be constant; as changing $e_i$ changes $\Lambda^2$ by $2\Lambda\cdot
S_i-4$, $\Lambda\cdot S_i=2$. Then $\Lambda\cdot F_l$, $\Lambda\cdot F_r$ are
even with difference $2$, the halves add nothing to
$\Theta=(\Lambda^2-\sigma)/4$, and a chain of negative-definite pieces
adds at most $1/2$, from its ends: the index must come from $b^+$. The
positive piece $P\cong\CP^2\#5\CPb$, $W_H$ capped by the adjacent traces,
has $\Theta(P)=1$ for the $\Lambda$ its chamber condition allows, and its
positive direction supplies $3/8$: $P$ raises $n_D$ by $5/8$. Its
Seiberg--Witten classes are controlled by a sphere of square zero and by toric
metrics with $\mathrm{Scal}\ge0$, admitting only $K$ with $K\cdot H$ between
$0$ and $\Lambda\cdot H$ ($H$ the period point); as $P$ is indefinite, a class
can pair nontrivially with both adjacent spheres at total cost $1/2$, not $1$.

\step{Step 6. The balance.}
With $n$ spheres and $m$ positive pieces, a copy of $B$ lowers $n_D$ by
$1/8$ and $\ell-T(v)$ by at least $3/8$; a positive piece raises $n_D$ by
$5/8$ and the bound on $\ell-T(v)$ by $7/8$, its $3/8$ plus the $1/2$ it saves
its neighbours. So $n_D=(5m-n)/8+O(1)$ and $\ell-T(v)\le(7m-3n)/8+O(1)$: for
large $n$, $n_D\ge1$ and exclusion hold if $1/5<m/n<3/7$, and outside
$[1/5,3/7]$ the count fails. As limits may break along the $3$-spheres, the inequality is needed on
each segment of consecutive pieces; it holds once positive pieces are three
copies of $B$ apart, except next to an irreducible component with nonzero
spinor, where the available argument needs four. So the average spacing is
in $[4,5)$; $u$ realizes four. Only $\phi$ allows it: $-W'$ has
$b^+=2$, so without $\phi$ the only return supplied by the triangles is $W_H$,
$m=n$, and exclusion fails. It must, since Step~2 also gives nonzero pairings
$\langle\Psi_+,(HB)^M\Psi_-\rangle$, with $n_D\to\infty$, for every
nontrivial knot. For $u$ and large $M$, the family
version of Step~1 gives $2^{n_D-1}\Omega_M=0$, while Step~2 gives
$\Omega_M\not\equiv0\pmod{2^N}$: the contradiction. That version is not in
the literature; Section~\ref{sec:closed} proves its closed model.

\section{A closed prototype}\label{sec:closed}

The mechanism of Steps~1 and~3 can be seen on a closed four-manifold, in a version of
\cite[Thm.~3.33(a)]{FL2b} in which the insertions exclude the reducibles,
though it cannot by itself give the contradiction. Let $X$ be closed,
oriented and simply connected with $b^+\ge1$, with generic metric $g$ and
perturbations as in \cite{FL2b,FLM}, and $\ft$, $w$, $\Lambda$, $\kappa$,
$n_D$ (FL's $n_a$) as in Step~1. For a spin$^c$ structure $\fs$ let
$v=c_1(\fs)-\Lambda\equiv w\pmod2$, and $M_\fs$ its Seiberg--Witten moduli
space, perturbed so that its points are the reducibles split by $\fs$
\cite[Lemma~3.13]{FL2a}; it occurs in the Uhlenbeck closure $\bar\cM_\ft$
only as
$M_\fs\times\Sym^\ell(X)$, $\ell=\kappa+v^2/4$
\cite[Lemma~3.32, (3.64)]{FL2b}. Let
$S_1,\dots,S_n$ be disjoint embedded spheres with $\langle w,S_i\rangle$ even,
and $T(v)=\#\{i:\langle v,S_i\rangle=0\}$.

\begin{theorem}\label{thm:closed}
Let $w\not\equiv0\pmod2$ and $n_D\ge1$, and let $z'$ be a monomial in the point
class and classes of embedded surfaces with $\deg z'+2n=\dim M^w_\kappa$.
\textup{(a)} If
\begin{equation}\label{eq:excl}
\kappa+v^2/4<T(v)\qquad\text{for every $\fs$ with $M_\fs\neq\emptyset$,}
\end{equation}
then $D^w_X(\mu(S_1)\cdots\mu(S_n)\,z')=0$ (for $b^+=1$, in the chamber of
$g$). \textup{(b)} Let $[S_i]=F_{l,i}-F_{r,i}$, where the $2n$ classes
$F_{l,i}$, $F_{r,i}$ are pairwise orthogonal of square $-2$ and $w$ is even on
each; let $\pi$ be the orthogonal projection to the complement of their span,
and $c_\perp$ the largest $(\pi v)^2$ for $M_\fs\neq\emptyset$. If
$2n>\deg z'+3(1+b^+)+2c_\perp$, then \eqref{eq:excl} holds.
\end{theorem}

For $n=0$, \eqref{eq:excl} is hypothesis (3.68) of \cite[Thm.~3.33(a)]{FL2b}
and (a) is that theorem. The insertions enter \eqref{eq:excl} on both sides:
each raises $\kappa$ by $1/4$ and forces a bubble onto every sphere
orthogonal to $v$.

\begin{proof}
(a) We follow the proof of \cite[Thm.~3.33(a)]{FL2b}, changing only the
representatives of $\mu(S_i)$; all new inputs are elementary except
transversality, standard but not written out for $\SO(3)$-monopoles. As $\langle w,S_i\rangle$ is even, $\bar\partial_A$ makes
$E_A|_{S_i}\otimes(\det E_A|_{S_i})^{-1/2}$ a holomorphic bundle $\mathcal
O(k)\oplus\mathcal O(-k)$ on $S_i\cong\mathbb P^1$ (Grothendieck), and the
coupled Dirac operator, $\bar\partial$ on its twist by $\mathcal O(-1)$, has
index zero and kernel $H^0(\mathcal O(k-1)\oplus\mathcal O(-k-1))$, nonzero
iff $k\ge1$. So the canonical section of its determinant line, of first Chern
class $\pm\mu(S_i)$ \cite{D90}, vanishes exactly on the jumping-line divisor
$V_i=\{[A]:k\ge1\}$. (J1)~$V_i$ depends only on $A|_{S_i}$ and is closed under
$C^0$ convergence of these restrictions, as $\bar\partial_A-\bar\partial_{A'}$
has order zero. (J2)~A connection preserving a splitting
$E|_{S_i}=L_1\oplus L_2$ with $\deg L_1=\deg L_2$ is not in $V_i$, as $k=0$.
A small perturbation through sections depending only on $A|_{S_i}$ gives
transversality on the irreducible strata of all levels, and makes $\cV(z)$
smooth and transverse to $M^w_\kappa$ at the finitely many points of
$\bar\cV(z)\cap\bar M^w_\kappa$, as \cite[Lemma~3.22]{FL2b} requires; it keeps (J1), and (J2) near the compact set of restrictions of
reducible limits with $\langle v,S_i\rangle=0$, where the canonical section
is nonzero. Put the
representatives of $z'$ and $\mu_c$ of \cite[Lemmas~3.12--3.13]{FL2b} in
general position, with base points off $\bigcup S_i$. Then no point lies in
supports of total codimension above four ($S_i$ counting two), i.e.\
\cite[(3.25)]{FL2b}: $z=\mu(S_1)\cdots\mu(S_n)\,z'$ is intersection-suitable.

So $\cZ=\cV(z)\cap\mathcal W^{n_D-1}\cap\cM^{*,0}_\ft/S^1$, where $\cV(z)$,
$\mathcal W$ represent $\mu_p(z)$, $\mu_c$ and $\cM^{*,0}_\ft$ consists of the irreducible monopoles
with nonzero spinor, is a one-manifold \cite[(3.26), (3.62)]{FL2b} with
$2^{n_D-1}D^w_X(z)$ signed ends at the link $\{\|\Phi\|^2_{L^2}=\varepsilon\}$
of the instantons: the proof of \cite[Prop.~3.29, (3.32)]{FL2b} uses only that
the representatives are determined by restriction, and the smoothness above.
Uhlenbeck convergence is $C^\infty$, after gauge, away from the bubbles
\cite[Def.~4.19, Thm.~4.20]{FL1}; so by (J1) and \cite[Lemma~3.15(1)]{FL2b} a
limit at level $\ell\ge1$ satisfies every condition whose support contains no
bubble, and as a level costs six dimensions and a bubble relieves at most
four, the count of \cite[Cor.~3.18]{FL2b} excludes irreducible limits with
nonzero spinor. There are no reducibles with zero spinor
\cite[Prop.~3.1, Cor.~3.3]{FL2a}. At a limit in
$M_\fs\times\Sym^\ell(X)$, a bubble lies on each $S_i$ with $\langle
v,S_i\rangle=0$, for otherwise the restrictions to $S_i$ would converge in $C^0$ to one
outside $V_i$, contradicting (J1) and (J2); the $S_i$ are disjoint, so
$\ell\ge T(v)$, contrary to \eqref{eq:excl}. Hence
$\bar\cZ\cap\{\|\Phi\|^2_{L^2}\ge\varepsilon\}$ is a compact oriented
one-manifold with boundary on the link, so $2^{n_D-1}D^w_X(z)=0$.

(b) The even numbers $\langle v,F_{l,i}\rangle$, $\langle
v,F_{r,i}\rangle$ are not both zero if $\langle v,S_i\rangle\neq0$, so
$-v^2/4=-(\pi v)^2/4+\sum_i(\langle v,F_{l,i}\rangle^2+\langle
v,F_{r,i}\rangle^2)/8\ge-c_\perp/4+(n-T(v))/2$; with
$8\kappa=2n+\deg z'+3(1+b^+)$ this gives \eqref{eq:excl}.
\end{proof}

In (b), $c_\perp$ depends on $\Lambda$. If $b^+\ge2$ and $K$ is basic, the
moduli spaces of $\pm K$ are nonempty, so $c_\perp\ge(\pi K)^2+(\pi\Lambda)^2$,
and (b) with $n_D\ge1$ needs
$b^-_\perp>4b^++7+\deg z'+(\pi K)^2-2n-\Lambda_R^2$, where $b^-_\perp$ counts
negative directions orthogonal to the halves and $\Lambda_R$ is the part of
$\Lambda$ along them. With $\Lambda$ constrained as in Step~5,
$\Lambda_R^2\le-2n$ and $n$ drops out (Step~4); with $\Lambda_R=0$ the halves
add $n/2$ to the Dirac index, which the sum over bundles forbids.

\begin{thebibliography}{DLME}
\bibitem[D]{D90} S.~K. Donaldson, \emph{Polynomial invariants for smooth
four-manifolds}, Topology \textbf{29} (1990), 257--315.
\bibitem[DLME]{DLME} A.~Daemi, T.~Lidman and M.~Miller Eismeier,
\emph{Filtered instanton homology and cosmetic surgery}, arXiv:2410.21248.
\bibitem[FL1]{FL1} P.~M.~N. Feehan and T.~G. Leness, \emph{PU(2) monopoles.
I}, J. Differential Geom. \textbf{49} (1998), 265--410.
\bibitem[FL2a]{FL2a} \bysame, \emph{PU(2) monopoles and links of top-level
Seiberg--Witten moduli spaces}, J. Reine Angew. Math. \textbf{538} (2001),
57--133.
\bibitem[FL2b]{FL2b} \bysame, \emph{PU(2) monopoles. II}, ibid., 135--212.
\bibitem[FL6]{FL6} \bysame, \emph{Witten's conjecture for many four-manifolds
of simple type}, J. Eur. Math. Soc. \textbf{17} (2015), 899--923.
\bibitem[FLM]{FLM} \bysame, \emph{An $\mathrm{SO}(3)$-monopole cobordism
formula relating Donaldson and Seiberg--Witten invariants}, Mem. Amer. Math.
Soc. \textbf{256} (2018), no.~1226.
\bibitem[KM]{KM04} P.~B. Kronheimer and T.~S. Mrowka, \emph{Witten's
conjecture and Property P}, Geom. Topol. \textbf{8} (2004), 295--310.
\bibitem[PT]{PT} V.~Ya. Pidstrigach and A.~N. Tyurin, \emph{Localisation of
Donaldson invariants along Seiberg--Witten classes}, arXiv:dg-ga/9507004.
\bibitem[W]{W} E.~Witten, \emph{Monopoles and four-manifolds}, Math. Res.
Lett. \textbf{1} (1994), 769--796.
\end{thebibliography}

\end{document}
